%% LyX 2.3.3 created this file.  For more info, see http://www.lyx.org/.
%% Do not edit unless you really know what you are doing.
\documentclass{article}
\usepackage[utf8]{inputenc}
\usepackage{geometry}
\usepackage{longtable}
\geometry{verbose,tmargin=1in,bmargin=1in,lmargin=1in,rmargin=1in}
\setlength{\parskip}{\medskipamount}
\setlength{\parindent}{0pt}
\usepackage{amsmath}
\usepackage{url}
\usepackage{amssymb}
\usepackage{setspace}
\onehalfspacing
\begin{document}
\begin{center}
{\Large{}Problem set \#3}{\Large\par}
\par\end{center}

\begin{center}
Due on November 25th by 12.30pm. Send solution to dealberti.francesco@gmail.com \\
Answers must be typed. You can work in teams under the constraint: Max(Team Size)=2 \\

\par\end{center}

\bigskip{}
\begin{enumerate}
\item (\textbf{Reweighting}) Go to David Autor's webpage and download
the cleaned 1979 and 1997 MORG files from:

https://economics.mit.edu/people/faculty/david-h-autor/data-archive

the paper is Trends in U.S. Wage Inequality: Revising the Revisionists.


a) Read through the cleaning programs associated with the MORG files.
Describe the important decisions being made.

b) Make a new variable ``lnhr\_wage'' which is the log of the ``hr\_wage''
variable. Restrict attention to the cleaned sample with ``hr\_wage\_sample==1''.
Use the ``kdensity'' command to plot kernel density estimates of
lnhr\_wage using the variable wgt\_hrs (which is the product of the
sample weights and the number of hours worked) as an ``aweight''.
Make density plots for both years on the same grid.

c) How has the wage distribution changed over time?

d) A number of demographic shifts occured between 1979 and 1997 that
might alter the aggregate wage distribution even if the distribution
for each demographic subgroup remained unaltered. Repeat the exercise
in b) separately for black men age 25-50 and for white women age 25-50.
How are the patterns different?

e) One way of adjusting the distribution for compositional differences
is to use propensity score reweighting.

Append the 1979 and 1997 files together and create a dummy variable
equal to one for observations in the 1997 sample. Estimate a logit
of the 1997 dummy on a dummy for female, a black dummy, an ``other
race'' dummy, and a quadratic polynomial in age. Use the ``aweights''
when estimating the logit.

f) Use the propensity score implied by this logit to reweight the
1997 wage distribution so that it has the 1979 demographic distribution.
Make a kernel density plot of the actual 1979 and 1997 wage distributions
against the ``counterfactual'' reweighted 1997 distribution.

g) Use the propensity score to reweight the 1979 wage distribution
so that it has the 1997 demographic distribution (if you get stuck
read the Dinardo, Fortin, and Lemieux paper). Make a kernel density
plot of the actual 1979 and 1997 wage distributions against the ``counterfactual''
reweighted 1979 distribution. Be sure to use the product of the wgt\_hrs
variable and the propensity score based weights when computing the
kernel density.

h) Based upon your analysis what would you say was the effect of the
demographic changes between 1979 and 1997 on the evolution of the
aggregate wage distribution?

i) Examine the suitability of your propensity score specification
by comparing the reweighted mean, standard deviation, and the 10th,
25th, 50th, 75th, and 90th percentiles of age in the 1997 sample to
the equivalent moments in the 1979 sample.

j) Do the same exercise restricting attention to blacks.

k) Try experimenting with different specifications for age in the
propensity score logit and checking how the results of i) and j) change.

l) What do you conclude about the suitability of the propensity score
specification?

\item (\textbf{Duflo, Hanna and Ryan (2012)}) This question asks you to solve and estimate a simplified version of the model from Duflo, Hanna and Ryan (2012). Consider a teacher deciding between work and leisure in two time periods, \(t \in \{0,1\}\). Let \(L_t \in \{0,1\}\) indicate whether the teacher takes leisure in period \(t\). The teacher's flow utility in period \(t\) is given by

\[
U_t = \beta C_t + L_t \cdot (\mu + \varepsilon_t),
\]

where \(C_t\) is consumption in period \(t\). All income and consumption is assumed to happen in period 2, so \(C_1 = 0\). In period 2, the teacher receives a salary payment of 5, and an extra bonus of 5 if she worked in both periods (that is, if \(L_1 = L_2 = 0\)). Assume that the teacher has a discount factor \(\delta \in [0,1]\).

\begin{enumerate}
    \item What are the state variables in period 2? Write an expression for the teacher's value function in period 2.
    
    \item Suppose that \(\varepsilon_t \sim \mathcal{N}(0,1)\). Compute the teacher's expected second-period payoff. Use this to write an expression for the teacher’s value function in period 1.
    
    \item Suppose you have a sample of teachers' choices \((L_{1i}, L_{2i})\) for \(i \in \{1, \dots, N\}\). Write an expression for the likelihood of teacher \(i\)'s choice sequence, \(\Pr[(L_1 = L_{1i}) \cap (L_2 = L_{2i})]\), as a function of the parameters \(\beta\), \(\mu\) and \(\delta\). Use this to write an expression for the sample log-likelihood.
    
    \item Download the data set \texttt{``dhr\_example.csv"} from the course website. This data set includes 800 observations on teachers' choice sequences \((L_{1i}, L_{2i})\). Import the data into a statistical software package. Estimate the parameters \((\beta, \mu, \delta)\) by maximum likelihood, and report your parameter estimates.
    
    \item Estimate an alternative version of the model where you assume that teachers are not forward-looking. Use your results from the two models to test the null hypothesis that teachers are not forward-looking. Discuss intuitively the features of the data that allow you to determine whether teachers are forward-looking.

    \item Estimate an alternative version of the model where you assume that teachers are not forward-looking. Use your results from the two models to test the null hypothesis that teachers are not forward-looking. Discuss intuitively the features of the data that allow you to determine whether teachers are forward-looking.

    \item Estimate an alternative version of the model where $\mu\sim N(\mu_{1};\sigma_{1})$. How does this version of the model differs from the previous one? Write down the associated likelihood and provide estimate of $(\beta,\delta, \mu_{1},\sigma_{1},$) along with their standard errors using SML. Make sure to explain how you derived the standard errors for your coefficients. 


    
    
\end{enumerate}
\end{enumerate}

\end{document}
